% =================================================
% Учительская версия рабочего листа
% =================================================

\worksheettitle
  {Точка, прямая, отрезок и луч}
  {Учительская версия: ответы и методические комментарии}

\begin{teacherbox}[Цель листа]
  Сформировать у пятиклассников устойчивое различение четырёх объектов:
  точки, прямой, отрезка и луча. Особое внимание уделяется школьной
  конвенции обозначений: прямая рисуется без стрелок, луч называется двумя
  буквами без знака вектора, первая буква указывает начало луча.
\end{teacherbox}

\section{Точка}

\pointsexamplefigure

\begin{propertybox}[Ответ]
  Точку изображают маленькой \textbf{меткой} и обозначают
  \textbf{заглавной} латинской буквой. На рисунке отмечены точки
  \(A\), \(B\), \(C\).
\end{propertybox}

\begin{teacherbox}[Диагностика]
  Проследите, чтобы буква располагалась рядом с меткой, но не заменяла её.
  Частая ошибка — поставить только букву без самой точки.
\end{teacherbox}

\section{Прямая}

\lineexamplefigure

\begin{propertybox}[Ответ]
  Прямая не имеет ни \textbf{начала}, ни \textbf{конца}. На рисунке показана
  только \textbf{часть} прямой.
\end{propertybox}

\begin{propertybox}[Через две точки]
  Через любые две различные точки можно провести прямую, и притом только
  одну.
\end{propertybox}

\incidenceexamplefigure

\begin{practicebox}[Ответы к заданию 2]
  \begin{enumerate}[label=\alph*),itemsep=0.35em]
    \item На прямой \(a\) лежат точки \(A\) и \(B\).
    \item Точка \(C\) не лежит на прямой \(a\).
    \item Прямая проходит через точки \(A\) и \(B\).
  \end{enumerate}
\end{practicebox}

\begin{teacherbox}[Точная речь]
  Полезно чередовать равносильные формулировки: «точка лежит на прямой» и
  «прямая проходит через точку». Пока можно не вводить символы
  \(A\in a\) и \(C\notin a\), если их ещё нет в используемом учебнике.
\end{teacherbox}

\section{Отрезок}

\segmentexamplefigure

\begin{propertybox}[Ответ]
  Точки \(A\) и \(B\) вместе со всеми точками прямой между ними образуют
  \textbf{отрезок} \(AB\). Точки \(A\) и \(B\) — \textbf{концы} отрезка.
  У отрезка два конца, и его длину можно измерить.
\end{propertybox}

\threelinepointsfigure

\begin{practicebox}[Ответ к заданию 3]
  Все отрезки с концами в отмеченных точках:
  \[
    AB,\qquad BC,\qquad AC.
  \]
\end{practicebox}

\begin{teacherbox}[Ключевая трудность]
  Некоторые ученики называют только «маленькие» отрезки \(AB\) и \(BC\),
  но не замечают отрезок \(AC\). Попросите провести пальцем путь от \(A\)
  до \(C\): точка \(B\) лежит внутри отрезка и не мешает существованию
  отрезка \(AC\).
\end{teacherbox}

\section{Луч}

\rayexamplefigure

\begin{propertybox}[Ответ]
  Луч имеет \textbf{начало}, но не имеет \textbf{конца}. Луч \(AB\)
  начинается в точке \(A\) и проходит через точку \(B\).
\end{propertybox}

\begin{notebox}[Обозначение]
  В тексте пишем «луч \(AB\)», без \(\overrightarrow{AB}\). Первая буква
  означает начало луча. На рисунке луч показан без стрелки: изображена лишь
  конечная часть объекта, который продолжается дальше.
\end{notebox}

\oppositeraysfigure

\begin{practicebox}[Ответы к заданию 4]
  \begin{enumerate}[label=\alph*),itemsep=0.35em]
    \item луч \(AB\);
    \item луч \(AC\);
    \item общее начало — точка \(A\).
  \end{enumerate}
\end{practicebox}

\begin{teacherbox}[Смысл порядка букв]
  Сопоставьте две фразы: «из \(A\) через \(B\)» и «из \(B\) через \(A\)».
  Это помогает понять, почему лучи \(AB\) и \(BA\) различаются без
  формального заучивания.
\end{teacherbox}

\section{Различаем объекты}

\identifyobjectsanswersfigure

\begin{center}
  \renewcommand{\arraystretch}{1.42}
  \begin{tabularx}{\textwidth}{
    >{\raggedright\arraybackslash}p{0.20\textwidth}
    >{\centering\arraybackslash}p{0.18\textwidth}
    >{\raggedright\arraybackslash}p{0.28\textwidth}
    >{\centering\arraybackslash}X}
    \toprule
    \rowcolor{TableHeader}
    Объект & Сколько концов? & Как продолжается? & Можно измерить длину? \\
    \midrule
    Прямая & нет & в обе стороны & нет \\
    Отрезок & два & не продолжается за концы & да \\
    Луч & один — начало & в одну сторону & нет \\
    \bottomrule
  \end{tabularx}
\end{center}

\section{Построение}

\constructionanswersfigure

\begin{teacherbox}[Проверка построения]
  Прямая \(AB\) должна продолжаться за обе точки. Отрезок \(CD\) заканчивается
  в точках \(C\) и \(D\). Луч \(DA\) начинается в \(D\), проходит через
  \(A\) и продолжается за \(A\). Все три объекта проводят сплошными линиями; различие определяется их
  концами и направлением продолжения.
\end{teacherbox}

\section{Верно или неверно?}

\begin{checkpointbox}[Ответы к заданию 6]
  \begin{enumerate}[label=\textbf{\arabic*.},itemsep=0.42em]
    \item Неверно: точки \(A\) и \(B\) лежат на прямой, но не являются её
      концами.
    \item Верно.
    \item Верно.
    \item Неверно: луч \(AB\) начинается в точке \(A\).
    \item Верно.
  \end{enumerate}
\end{checkpointbox}

\begin{teacherbox}[Ожидаемый итоговый ответ]
  Прямая не имеет ни начала, ни конца. Отрезок имеет два конца. Луч имеет
  начало, но не имеет конца.
\end{teacherbox}

\section{Типичные ошибки}

\begin{center}
  \renewcommand{\arraystretch}{1.4}
  \begin{tabularx}{\textwidth}{
    >{\raggedright\arraybackslash}p{0.31\textwidth}
    >{\raggedright\arraybackslash}p{0.30\textwidth}
    >{\raggedright\arraybackslash}X}
    \toprule
    \rowcolor{TableHeader}
    Ошибка & Что не понято & Наводящий вопрос \\
    \midrule
    У прямой поставлены стрелки
      & Не усвоена принятая школьная модель рисунка
      & Где на рисунке находятся концы прямой? \\
    \addlinespace[0.2em]
    Луч назван \(BA\), хотя начинается в \(A\)
      & Не понято значение первой буквы
      & Из какой точки выходит луч? \\
    \addlinespace[0.2em]
    Отрезки \(AB\) и \(BA\) считаются разными
      & Порядку букв ошибочно приписано направление
      & Изменились ли концы отрезка? \\
    \addlinespace[0.2em]
    На трёх точках найдено только два отрезка
      & Не замечен составной отрезок
      & Есть ли отрезок с концами \(A\) и \(C\)? \\
    \bottomrule
  \end{tabularx}
\end{center}
